% Lösungen Einheit 1 von 3 – Potenzen (zusammenbau v0.1)
\einheitenkopf{Einheit 1 von 3 · Potenzen}

\erg{12}{a) positiv – das Minus steht in der Klammer, die Hochzahl ist gerade}
\erg{13}{a) Bruch – das Minus steht in der Hochzahl: eins geteilt durch die Potenz}
\erg{14}{a) hoch – eine Potenz aus vier gleichen Faktoren \quad b) $2 \cdot 2 \cdot 2 \cdot 2 \cdot 2 \cdot 2 = 64$ \quad c) $7^2 = 49$ und $7 \cdot 2 = 14$ \quad d) $196$ \quad e) $529$ \quad f) $(-4) \cdot (-4) \cdot (-4) = -64$ \quad g) $-(8 \cdot 8) = -64$ \quad h) $0{,}6 \cdot 0{,}6 = 0{,}36$ \quad i) $2^7 = 128$ ist größer als $5^3 = 125$ \quad j) $5^3 = 125$, $2^6 = 64$, $9^2 = 81$ – die größte ist $5^3$ \quad k) $5 \cdot 5 \cdot 5 \cdot 5 = 625$, also $x = 4$ \quad l) $4\,096 = 64 \cdot 64$ und $64 = 4 \cdot 4 \cdot 4$, also $x = 6$ \quad m) vier Nullen, also $x = 4$ \quad n) $4^{-2} = \frac{1}{4^2} = \frac{1}{16} = 0{,}0625$ \quad o) $4^{-3} = \frac{1}{64} = 0{,}015625 < 0{,}02$ \quad p) $2, 4, 8, 16, 32, 64, 128, 256, 512$: neun Faktoren, also $x = 9$}
\erg{15}{a) $\frac{2 \cdot 2}{5 \cdot 5} = \frac{4}{25}$}
\erg{16}{a) $0{,}5^2 = 0{,}25 > 0{,}2 = 20\,\%$}
\erg{17}{a) $1$}
\erg{18}{a) Basis mal Exponent gerechnet; richtig: $6^3 = 6 \cdot 6 \cdot 6 = 216$ \quad b) Man muss ausrechnen, die größere Basis entscheidet nicht: $3^5 = 243$, aber $5^3 = 125$.}
